% TOP_PROBLEM: {"id": 30001499, "title": "Novikov Homotopy Invariance Conjecture", "category_id": 7, "category": {"id": 7, "name": "topology", "display_name": "Topology", "description": "Properties preserved under continuous deformations.", "slug": "topology", "order_index": 7, "created_at": "2026-07-31T15:26:25.671Z"}, "status": "open"}
% FIRST_POSTED: 2026-09-27
% !TeX program = lualatex
% Compile twice: lualatex 30001499.tex
% Original entry retained; research subsection and [E] references added.
% No external bibliography, image, data, or style file is required.
\documentclass[11pt,a4paper]{article}
\usepackage[margin=24mm,headheight=14pt]{geometry}
\usepackage{fontspec}
\setmainfont{Latin Modern Roman}
\setsansfont{Latin Modern Sans}
\setmonofont{Latin Modern Mono}
\usepackage{amsmath,amssymb,mathtools,amsthm}
\usepackage{unicode-math}
\setmathfont{Latin Modern Math}
\usepackage{microtype}
\usepackage{enumitem,xurl,needspace}
\usepackage[dvipsnames]{xcolor}
\usepackage{hyperref}
\hypersetup{unicode=true,colorlinks=true,linkcolor=MidnightBlue,
 urlcolor=MidnightBlue,bookmarksnumbered=true,
 pdftitle={Research notebook - Problem 30001499},
 pdfauthor={Research attempt; not independently refereed}}
\usepackage{bookmark}
\usepackage{titlesec}
\titleformat{\section}{\Large\bfseries\raggedright}{\thesection}{0.7em}{}
\usepackage{fancyhdr}
\pagestyle{fancy}\fancyhf{}
\fancyhead[L]{\small\sffamily Research notebook}
\fancyhead[R]{\small\sffamily Problem 30001499}
\fancyfoot[L]{\scriptsize 27 September 2026}
\fancyfoot[C]{\thepage}
\fancyfoot[R]{\scriptsize Unrefereed research attempt}
\setlength{\parindent}{0pt}
\setlength{\parskip}{5pt plus 1pt}
\setlength{\emergencystretch}{3em}
\clubpenalty=10000
\widowpenalty=10000
\displaywidowpenalty=10000
\setcounter{secnumdepth}{1}
\setlist[itemize]{leftmargin=3em,itemsep=5pt,topsep=4pt,parsep=0pt}
\allowdisplaybreaks
\newcommand{\N}{\mathbb{N}}
\newcommand{\Z}{\mathbb{Z}}
\newcommand{\Q}{\mathbb{Q}}
\newcommand{\R}{\mathbb{R}}
\newcommand{\C}{\mathbb{C}}
\newcommand{\F}{\mathbb{F}}
\newcommand{\A}{\mathbb{A}}
\newcommand{\PP}{\mathbb P}
\newcommand{\E}{\mathbb{E}}
\newcommand{\eps}{\varepsilon}
\newcommand{\dd}{\,\mathrm d}
\newcommand{\sref}[2]{\hyperlink{source:#1:#2}{[S#2]}}
\newcommand{\eref}[2]{\hyperlink{extra:#1:#2}{[E#2]}}
\newif\ifshowcatalogue
\showcataloguetrue
\newcommand{\cataloguescope}[1]{\ifshowcatalogue
 \paragraph{Exact scope in the supplied catalogue.}
 \begin{quote}\small #1\end{quote}\fi}
\newtheorem{notebooklemma}{Lemma}[section]
\newtheorem{notebookproposition}[notebooklemma]{Proposition}
\numberwithin{equation}{section}
\begin{document}
\setcounter{section}{0}
% ================================================================
% problemId: 30001499
% Canonical problem ID: 30001499
% exactTarget SHA-256: ae2925817907faa2e3585a66215a53756b692dde79bc142a05f2219493b8ef6c
\clearpage
\section*{\texorpdfstring{Novikov Homotopy Invariance Conjecture}{Novikov Homotopy Invariance Conjecture}}\label{30001499}
\subsection{Definitions and mathematical statement}
Let $\pi=\pi_1(M)$ and let $B\pi$ be a classifying space with contractible universal cover. A map $r:M\to B\pi$ classifies the universal covering. The Hirzebruch class $L(M)\in\bigoplus_jH^{4j}(M,\Q)$ is the characteristic class built from the Pontryagin classes of $TM$. For $u\in H^*(B\pi,\Q)$, its higher signature is
\[
 \sigma_u(M,r)=\langle L(M)\smile r^*u,[M]\rangle,
\]
where only terms of total degree $\dim M$ contribute. For every orientation-preserving homotopy equivalence $f:N\to M$, require $\sigma_u(N,r\circ f)=\sigma_u(M,r)$.
\par\smallskip\textit{Formulation and concept references:} \sref{30001499}{1}.
\cataloguescope{For every closed oriented smooth manifold M, every u in H\textasciicircum{}*(B pi\_1(M); Q), and every classifying map r, prove that the higher signature <L(M) cup r\textasciicircum{}*u,[M]> is invariant under orientation-preserving homotopy equivalence.}
\subsection{Short English statement}
Are all higher signatures unchanged by orientation-preserving homotopy equivalence of closed manifolds?
\subsection{Research attempt}

\paragraph{Current frontier.}
The operator-theoretic sufficient condition is rational injectivity of the assembly map from $K_*(B\Gamma)$; the signature index is homotopy invariant, and injectivity recovers all higher signatures. Leichtnam--Piazza give this implication in Theorem 7.10, with the normalization in equation (7.9) \eref{30001499}{1}. Oyono-Oyono--Yu develop quantitative index methods and geometric proofs for groups with finite decomposition complexity; this is a structural hypothesis, not a theorem for every group \eref{30001499}{2}. Tian--Yu's September 2026 revision states a further sufficient condition: a finite-complexity coarse embedding into a universal Banach space implies the strong Novikov conjecture \eref{30001499}{3}. That recent preprint is used below as a stated theorem, not as an independently audited result. The elementary embedding calculations in this attempt do not depend on its proof.

\paragraph{Attack route.}
One could seek a universal inverse for the rational assembly map, or extend quantitative decomposition arguments to every group; both require substantial new index-theoretic input. The route tried here is different: all bounded-degree graphs admit a particularly elementary coarse embedding into
\[
 U=\left(\bigoplus_{p\geq1}\ell^{2p}(\N)\right)_{\ell^2}.
\]
Can reweighting its coordinates impose the extra control required by the recent theorem? We first prove a subsequence criterion for success, then prove that scalar reweighting of the natural universal construction necessarily fails. This identifies the precise improvement that a different family of maps would need.

\paragraph{Attempt: the quantitative input.}
For real sequences let
\[
 \phi_p:S(\ell^{2p})\longrightarrow S(\ell^2),\qquad
 \phi_p(x)_i=\operatorname{sgn}(x_i)|x_i|^p.
\]
These are the standard Mazur maps used in \eref{30001499}{3}, Definition 1.2. Their Lipschitz constant is $p$. For the upper bound, along the segment between two unit vectors, H\"older's inequality gives
$\|D\phi_p(x)v\|_2\leq p\|x\|_{2p}^{p-1}\|v\|_{2p}\leq p\|v\|_{2p}$; integrate along the segment in the unit ball. For the reverse inequality take
$x=(2^{-1/(2p)},2^{-1/(2p)})$ and the unit tangent vector
$v=(2^{-1/(2p)},-2^{-1/(2p)})$. A normalized sphere path with this tangent has image derivative of $\ell^2$-norm $p$.

The part of finite complexity we will verify explicitly is the existence of coordinate bounds
\begin{equation}\label{r56:bounds}
 \rho_p^-(d(x,y))\leq\|h_p(x)-h_p(y)\|_{2p}
                      \leq\rho_p^+(d(x,y))
\end{equation}
with
\begin{equation}\label{r56:complexity}
 \sum_p\rho_p^-(t)^2\longrightarrow\infty\quad(t\to\infty),
 \qquad
 \sum_p(p+1)^2\rho_p^+(t)^2<\infty\quad(t<\infty).
\end{equation}
We impose the upper condition on the whole sum, which is stronger than allowing finitely many exceptional coordinates at each scale as in the source. We also check directly that $h(x)\in U$. Thus the construction below supplies, rather than assumes, the relevant analytic hypotheses.

Let $\Gamma$ be an infinite finitely generated group with its word metric. It is countable, so $\ell^{2p}(\Gamma)$ can be identified isometrically with a coordinate subspace of $\ell^{2p}(\N)$.

\begin{notebookproposition}[Subsequence improvement criterion]\label{r56:selection}
Suppose there are distinct increasing integers $p_j$ and maps
$\tau_j:\Gamma\to\ell^{2p_j}(\Gamma)$ with Lipschitz constants $L_j>0$. Suppose also that some $s_j>0$ and $R_j<\infty$ satisfy
\[
 d(x,y)\geq R_j\quad\Longrightarrow\quad
       \|\tau_j(x)-\tau_j(y)\|_{2p_j}\geq s_j.
\]
If
\begin{equation}\label{r56:missing}
 \liminf_{j\to\infty}\frac{(p_j+1)L_j}{s_j}=0,
\end{equation}
then a subsequence of these maps, after rescaling and centering, defines a coarse embedding satisfying \eqref{r56:bounds}--\eqref{r56:complexity}.
\end{notebookproposition}
\begin{proof}
Choose indices $j_k$ such that
$((p_{j_k}+1)L_{j_k}/s_{j_k})^2\leq2^{-k}$. We may enlarge $R_{j_k}$ to at least $k$, without changing its separation property. For a fixed identity element $e$ put
\[
 h_{p_{j_k}}(x)=\frac{\tau_{j_k}(x)-\tau_{j_k}(e)}{s_{j_k}},
 \qquad h_p(x)=0\ \text{on other coordinates}.
\]
For fixed $x$, the sum of the squared norms is at most
$d(x,e)^2\sum_k(L_{j_k}/s_{j_k})^2<\infty$.
Take $\rho_{p_{j_k}}^+(t)=tL_{j_k}/s_{j_k}$ and
$\rho_{p_{j_k}}^-(t)=\mathbf1_{[R_{j_k},\infty)}(t)$, with zero bounds on inactive coordinates. The weighted upper sum is at most $t^2\sum_k2^{-k}$. For every $K$, the lower sum is at least $K$ when $t\geq\max_{k\leq K}R_{j_k}$. Enlarging the radii to at least $k$ also makes this lower sum finite for each fixed $t$. This proves all assertions, including properness.
\end{proof}

There is a useful exact explanation for \eqref{r56:missing}. In a scalar weighting scheme write
$v_j=w_j^2s_j^2$ and $a_j=((p_j+1)L_j/s_j)^2$. To make the lower contributions diverge while keeping the corresponding upper estimate summable, one asks for
\begin{equation}\label{r56:weights}
 \sum_jv_j=\infty,\qquad \sum_ja_jv_j<\infty.
\end{equation}
For positive $a_j$, such nonnegative finite weights exist if and only if
$\liminf a_j=0$. The subsequence argument proves sufficiency; if eventually $a_j\geq c>0$, the second sum forces the tail of the first to be finite. This equivalence concerns this particular lower/upper-bound bookkeeping. It is not a characterization of all possible embeddings.

\paragraph{Attempt: testing the universal radial construction.}
Let the Cayley graph have degree at most $D\geq2$. For integer $p\geq1$ define
\[
 \theta_p(x)(z)=\left(1-\frac{d(x,z)}p\right)_+,
 \qquad z\in\Gamma.
\]
The support is contained in $B(x,p-1)$, which has at most $D^p$ vertices. Consequently
\begin{equation}\label{r56:radialbounds}
 \|\theta_p(x)\|_{2p}\leq\sqrt D,\qquad
 \frac1p\leq\operatorname{Lip}(\theta_p)
                   \leq\frac{\sqrt{2D}}p.
\end{equation}
For the upper Lipschitz bound, adjacent centers change each coordinate by at most $1/p$, and the union of supports has at most $2D^p$ elements. Extend the edge estimate along paths. For the lower bound, at coordinate $z=x$ adjacent centers have values $1$ and $1-1/p$. When $d(x,y)\geq2p$, supports are disjoint and a center coordinate gives
\begin{equation}\label{r56:separation}
 \|\theta_p(x)-\theta_p(y)\|_{2p}\geq1.
\end{equation}

The centered direct sum
\[
 F(x)=\left(p^{-1/2}(\theta_p(x)-\theta_p(e))\right)_{p\geq1}
\]
is a coarse embedding into $U$. The squared upper bound is at most
$2D\,d(x,y)^2\sum_pp^{-3}$, so the map is defined and Lipschitz. The squared lower bound is at least
$\sum_{p\leq d(x,y)/2}1/p$, which tends to infinity. This proof holds for every bounded-degree Cayley graph, without amenability, a growth assumption, or a torsion-free assumption.

It is tempting to alter the weights and then apply the finite-complexity theorem. The next calculation rules out that entire scalar-weighting repair, for these coordinates and these Mazur maps.

\begin{notebookproposition}[A scalar-reweighting obstruction]\label{r56:obstruction}
For any real scalars $w_p$, set
$h_p(x)=w_p(\theta_p(x)-\theta_p(e))$. If this map satisfies the weighted upper finite-complexity condition at distance $1$, even after omitting finitely many coordinates, then $\sum_pw_p^2<\infty$ and the image of $h$ is bounded. In particular it cannot be a coarse embedding of the infinite group.
\end{notebookproposition}
\begin{proof}
For an edge $(x,y)$, evaluating coordinate $z=x$ shows that every valid upper function satisfies $\rho_p^+(1)\geq|w_p|/p$. Since $\operatorname{Lip}(\phi_p)=p$,
\[
 (p+1)^2\rho_p^+(1)^2\geq w_p^2.
\]
Summability of the tail forces square summability of all the weights. But \eqref{r56:radialbounds} then gives the uniform bound
\[
 \|h(x)-h(y)\|_U^2\leq4D\sum_pw_p^2<\infty.
\]
The word metric is unbounded, contradicting the required lower control.
\end{proof}

\paragraph{Attempt: changing the radius, not merely the weights.}
One possible escape is to use radius $R$ in $\ell^{2p}$ instead of radius $p$. If $V(R)=|B(e,\lceil R\rceil)|$, the same proof gives a separation lower bound $1$ beyond $2R$ and an edge upper bound
\begin{equation}\label{r56:larger-radius}
 L\leq\frac{(2V(R+1))^{1/(2p)}}{R}.
\end{equation}
This really helps for subexponential growth. Put
$b(R)=\max\{1,\log(2V(R+1))\}=o(R)$, and choose
$p=\lceil\sqrt{Rb(R)}\rceil$. Along a subsequence with distinct $p$, both $p/R\to0$ and $b(R)/p\to0$, so the right side of $(p+1)L$ tends to zero. Proposition~\ref{r56:selection} applies. This is a consistency check in a familiar favorable class, not a new resolution for that class.

For a merely exponential growth estimate $V(R)\leq C\exp(hR)$, optimizing the available upper bound produces $R$ proportional to $p$ and no decay to zero in $pL$. That observation alone does not prove impossibility for other maps: an upper bound that fails to improve is not a lower-bound theorem. The proved impossibility is only Proposition~\ref{r56:obstruction} for the specified radial blocks.

\paragraph{Stress test and connection back to the exact target.}
If maps satisfying \eqref{r56:missing} were constructed for every infinite finitely generated group, and the stated Tian--Yu theorem is valid in its inspected revision, Proposition~\ref{r56:selection} would give rational assembly injectivity for every such group \eref{30001499}{3}. The homotopy invariance of the signature index would then give
\[
 r_*\bigl(L(M)\cap[M]\bigr)
  =(r\circ f)_*\bigl(L(N)\cap[N]\bigr)
       \quad\text{in }H_*(B\pi;\Q),
\]
and pairing with every $u$ is precisely the original conclusion \eref{30001499}{1}. Fundamental groups of closed connected manifolds are finitely generated. Finite fundamental groups contribute no positive-degree rational group cohomology, and the remaining ordinary signature is already homotopy invariant.

Thus the missing improvement would establish a strong Novikov assertion, not merely a special case of the displayed higher-signature identity. It would \emph{not}, by this reasoning, prove the surjectivity or the all-locally-compact-group statement of catalogue Section 53. The distinction between $B\pi$ here and the proper-action classifying space in Baum--Connes is retained.

No uncontrolled interchange of infinitely many coordinates is needed: all upper sums have explicit summable majorants, and properness follows by activating a finite number of coordinates at a time. The accompanying script checks the radial edge and separation estimates on the integer line for $1\leq p\leq20$ with exact rational arithmetic; this is only a check on the formulas. No finite group or finite graph computation could establish the required universal map improvement.

\paragraph{Outcome.}
\textbf{Outcome: Conditional reduction.}

The subsequence criterion and the scalar-reweighting obstruction are proved for the specified maps and control functions. Together they identify a concrete target: an unbounded sequence of blocks with $(p+1)L/s$ tending to zero along a subsequence. The universal radial construction has the wrong scale and cannot be repaired by scalar weights alone. No such improvement is obtained for arbitrary groups, and the final index-theoretic implication additionally relies on the stated recent preprint theorem. No complete Novikov proof is claimed, and novelty of the intermediate deductions has not been established by an independent literature audit.

\subsection{Sources}
\begin{itemize}\raggedright
\item[\textnormal{[S1]}]\hypertarget{source:30001499:1}{} Eric Leichtnam and Paolo Piazza, Elliptic Operators and Higher Signatures, Question 1 (2004).
\url{https://aif.centre-mersenne.org/item/10.5802/aif.2049.pdf}.
{\small\textit{Catalogue formulation source.}}
\item[\textnormal{[S2]}]\hypertarget{source:30001499:2}{} Novikov's Conjecture.
\url{https://arxiv.org/abs/1506.05408}.
{\small\textit{Catalogue research/status reference; not a proof endorsement.}}
\item[\textnormal{[S3]}]\hypertarget{source:30001499:3}{} Quantitative index, Novikov conjecture and coarse decomposability.
\url{https://arxiv.org/abs/2412.01314}.
{\small\textit{Catalogue research/status reference; not a proof endorsement.}}
\item[\textnormal{[S4]}]\hypertarget{source:30001499:4}{} Embedding complexity into the universal Banach space and the strong Novikov conjecture.
\url{https://arxiv.org/abs/2605.12930}.
{\small\textit{Catalogue research/status reference; not a proof endorsement.}}

% Research references added for this notebook.
\item[\textnormal{[E1]}]\hypertarget{extra:30001499:1}{} Eric Leichtnam and Paolo Piazza, \emph{Elliptic Operators and Higher Signatures}, arXiv:math/0406020 (2004), Section 7.6, equation (7.9), Theorem 7.10, and Section 7.7.
\url{https://arxiv.org/pdf/math/0406020}.
{\small\textit{Accessible author manuscript of the formulation reference; rational assembly injectivity and the signature-index argument. Consulted 27 September 2026.}}
\item[\textnormal{[E2]}]\hypertarget{extra:30001499:2}{} Herv\'e Oyono-Oyono and Guoliang Yu, \emph{Quantitative index, Novikov conjecture and coarse decomposability}, arXiv:2412.01314v1 (2 December 2024), Introduction and the quantitative index/decomposition results.
\url{https://arxiv.org/html/2412.01314v1}.
{\small\textit{Finite decomposition complexity is a hypothesis, not a universal property of groups. Consulted 27 September 2026.}}
\item[\textnormal{[E3]}]\hypertarget{extra:30001499:3}{} Geng Tian and Guoliang Yu, \emph{Embedding complexity into the universal Banach space and the strong Novikov conjecture}, arXiv:2605.12930v7 (15 September 2026), Theorem 1.1, Definition 1.2, and the discussion on pages 2--4.
\url{https://arxiv.org/pdf/2605.12930v7}.
{\small\textit{Recent preprint theorem and exact finite-complexity conditions. The full index-theoretic proof has not been independently audited here. Consulted 27 September 2026.}}
\end{itemize}

\end{document}
